\documentclass[11pt]{article}

\usepackage[margin=1in]{geometry}
\usepackage{amsmath, amssymb, amsthm}
\usepackage{mathtools}
\usepackage{booktabs}
\usepackage{graphicx}
\graphicspath{{figures/}{../figures/}}
\usepackage[hidelinks]{hyperref}
\usepackage[authoryear,round]{natbib}
\bibpunct{(}{)}{,}{a}{}{,}

\newcommand{\R}{\mathbb{R}}
\newcommand{\E}{\mathbb{E}}
\newcommand{\F}{\mathcal{F}}

\theoremstyle{plain}
\newtheorem{theorem}{Theorem}
\newtheorem{lemma}{Lemma}
\newtheorem{proposition}{Proposition}

\title{Optimal Trading Filters:\\ a Wiener--Hopf Approach}
\author{Matteo Smerlak\thanks{Capital Fund Management, 23 rue de l'Universit\'e, 75007 Paris, France.\ \texttt{matteo.smerlak@cfm.com}}}
\date{}

\begin{document}
\maketitle

\begin{center}
\begin{tabular}{p{0.9\linewidth}}
\textbf{Keywords.} Optimal trading; market impact; propagator model; inventory risk; Wiener--Hopf factorization; linear filters; fractional calculus; adapted stochastic control.\\[2pt]
\textbf{JEL classification.} C61; G11; G12; G14.\\[2pt]
\textbf{MSC classification.} 45E10; 60G35; 91G80; 26A33.
\end{tabular}
\end{center}

\begin{abstract}
We solve the quadratic gain--risk--cost trading problem in which a trader monetizes a return-predicting signal against three frictions: temporary impact, transient impact with a propagator kernel, and mean--variance inventory risk. The friction operator $Q$ is non-local in time, so the unconstrained optimum needs future signal values and adaptedness is the binding constraint. A Wiener--Hopf factorization $Q=Q_-Q_+$ into causal and anticausal factors resolves it: the adapted optimum is $u^\star = Q_+^{-1}P_+Q_-^{-1}\alpha$, with $P_+$ the projection onto adapted processes. For stationary Gaussian signals this is a linear trading filter. It reduces to the Markowitz rule under risk alone and to a fractional derivative of order $(1-\beta)/2$ under power-law impact alone; rational symbols give filters with finitely many exponential moving averages, recovering the G\^arleanu--Pedersen aim portfolio and the stationary Neuman--Vo\ss{} solution. The causal factor on the imaginary axis yields a closed-form rate response and a phase diagram for when the optimum trades against its own transient impact. On a finite horizon the factors become terminal-anchored Gohberg--Krein factors, recovering the Gatheral--Schied--Slynko and Forde--S\'anchez-Betancourt--Smith solutions. All formulas are checked against discretized adapted optima.
\end{abstract}

% =====================================================
\section{Introduction}

\subsection{Gain--risk--cost: a quadratic optimization}

A trader holds a return-predicting signal $\alpha_t$ --- the conditional expectation, at time $t$, of the asset's remaining price appreciation --- and chooses an adapted trading rate $u_t$, accumulating position $x_t = \int_{-\infty}^t u_s\,ds$. Trading is costly in three ways: a temporary impact cost quadratic in the instantaneous rate; a transient impact cost in which each trade moves the price against later trades, with an influence decaying according to a propagator kernel $G$ \citep{BouchaudGefenPottersWyart2004,Gatheral2010}; and an inventory penalty proportional to the variance of the position. The problem is
\begin{equation}\label{eq:objective}
\max_{u\ \text{adapted}}\ \E\!\int u_t\alpha_t\,dt
\;-\; \frac{\eta}{2}\,\E\!\int u_t^2\,dt
\;-\; \frac{\gamma}{2}\,\E\!\iint G(|t-s|)\,u_t u_s\,dt\,ds
\;-\; \frac{\lambda}{2}\,\E\!\int x_t^2\,dt .
\end{equation}
Empirical propagators decay as power laws $G(t) = |t|^{-\beta}$ with $\beta\in(0.2,0.6)$ across markets \citep{BouchaudGefenPottersWyart2004,JusselinRosenbaum2020}; the exponential kernel $G(t)=e^{-\kappa|t|}$ underlies the resilience model of \citet{ObizhaevaWang2013} and recovers the temporary and permanent costs of \citet{AlmgrenChriss2001} in its fast- and slow-decay limits.

Problem \eqref{eq:objective} is a linear gain against a positive quadratic form. On the whole line the form is diagonal in frequency, with symbol
\begin{equation}\label{eq:symbol}
Q(\omega) \;=\; \eta \;+\; \gamma\,\hat G(\omega) \;+\; \frac{\lambda}{\omega^2},
\qquad \hat G(\omega) = \int e^{i\omega t}G(|t|)\,dt \ge 0,
\end{equation}
where the risk term acts on the position, $|\hat x|^2 = |\hat u|^2/\omega^2$. Were information no constraint, the solution would be one line: $u = Q^{-1}\alpha$, where $Q$ is the operator with symbol $Q(\omega)$ (convention of \S2.1), the continuous-time analog of the Markowitz rule $w = \lambda^{-1}\Sigma^{-1}\mu$ (weights $w$, covariance $\Sigma$, expected returns $\mu$), with the friction symbol in place of the covariance matrix.

\subsection{Adaptedness: the key constraint when costs are non-local}

With temporary cost alone ($\gamma=\lambda=0$) the Hessian is diagonal in time: the first-order condition at time $t$ involves $u_t$ and $\alpha_t$ only, and the adapted optimum is the myopic rule $u_t = \alpha_t/\eta$: conditional expectation passes through a pointwise equation at no cost. Transient impact and inventory risk destroy this locality. The operator $Q$ couples different times: $(Qu)(t)$ integrates the trading rate over the past \emph{and the future}, so the unconstrained optimum $(Q^{-1}\alpha)_t$ is a functional of signal values the trader has not yet observed. Adaptedness --- $u_t$ measurable with respect to $\F_t$ --- becomes the binding constraint, and it is a constraint of an unusual kind: a closed subspace $L^2_{\rm adap}$ cut out by an increasing family of $\sigma$-algebras. Varying \eqref{eq:objective} over adapted perturbations gives the \emph{projected} first-order condition
\begin{equation}\label{eq:foc}
\E_t\bigl[(Qu^\star)(t)\bigr] \;=\; \alpha_t
\qquad\Longleftrightarrow\qquad
P_+\,Q\,P_+\,u^\star = \alpha,
\end{equation}
with $P_+$ the orthogonal projection onto adapted processes, $(P_+X)_s = \E_s[X_s]$. Solving \eqref{eq:foc} requires inverting $P_+QP_+$; the inverse of a projected operator is not the projection of the inverse, precisely because $Q$ is non-local.

\subsection{The Wiener--Hopf method}

\citet{WienerHopf1931} introduced the technique that resolves this class of problems, originally for convolution equations on a half-line, $\int_0^\infty k(t-s)f(s)\,ds = g(t)$ for $t\ge0$ \citep{Krein1962,Noble1958}. Fourier analysis diagonalizes convolutions on the whole line; a half-line support constraint couples the frequencies, in the same way the filtration constraint couples them in \eqref{eq:foc}. The method restores solvability by factoring the symbol,
\[
\hat k(\omega) \;=\; k_-(\omega)\,k_+(\omega),
\]
with $k_+$ analytic and zero-free in the upper half-plane and $k_-$ in the lower. Analyticity in a half-plane encodes causality: by Paley--Wiener, $k_+$ is the transform of a kernel supported on $t\ge0$ and $k_-$ of one supported on $t\le0$. Operators with causal kernels map functions supported on a half-line $[a,\infty)$ into functions supported on the same half-line: they are triangular with respect to the time order, and triangular operators invert forward in time, staying triangular. The factorization thus splits a two-sided convolution into two one-sided operators, each compatible with the boundary, between which the half-line projection can be threaded. The finite-dimensional analog of the construction is the Cholesky factorization of a positive matrix into triangular factors, computed by Gram--Schmidt along the coordinate order.

The method's classical applications are prediction and control. For linear prediction of a stationary time series, the factored symbol is the spectral density, $S = |\varphi_+|^2$, and the causal factor builds the innovations representation \citep{Wiener1949,Whittle1963}. In the trading problem the factored symbol is the \emph{friction} $Q$, and the causal factor builds the optimal trade.

\subsection{This paper}

We solve \eqref{eq:foc} with a stochastic-process version of the Wiener--Hopf method, in which the half-line projection of the classical theory is replaced by the optional projection $P_+$ onto adapted processes. The algebraic core is the projected-inverse identity
\begin{equation}\label{eq:pi-intro}
(P_+\,Q\,P_+)^{-1} \;=\; Q_+^{-1}\,P_+\,Q_-^{-1}
\end{equation}
for any factorization $q = Q_-Q_+$ into a causal factor with causal inverse and its anticausal adjoint (Lemma~\ref{lem:pi}); the abstract setting is Arveson's factorization of positive operators on the nest (increasing family) of adapted subspaces \citep{Arveson1975}. The resulting policy is explicit for \emph{any} adapted signal: anticausally whiten the trader's forecast curve, project onto the filtration, causally color (Theorem~\ref{thm:general}). For stationary Gaussian signals carrying their own filtration the policy becomes a linear filter with computable transfer function (Theorem~\ref{thm:filter}), and the paper's remaining sections read off its consequences: the Markowitz and fractional-derivative limits (\S4); explicit filters that recover the aim-portfolio rule of \citet{GarleanuPedersen2013,GarleanuPedersen2016} and make explicit, in the stationary regime, the FBSDE solutions of \citet{NeumanVoss2022} (\S5); a closed-form rate response that identifies when the optimum trades against the signal to ride its own transient impact (\S6); and the finite-horizon theory, in which the Wiener--Hopf factors become terminal-anchored Gohberg--Krein factors and the Fredholm solutions of \citet{FordeSanchezSmith2022} and the U-shaped liquidation profiles of \citet{GatheralSchiedSlynko2012} are recovered (\S7). All closed forms are verified against discretized adapted optima (\S6.3).

Relative to the execution literature --- deterministic liquidation under general kernels \citep{GatheralSchiedSlynko2012}, signal-adaptive trading under temporary and exponential-resilience costs \citep{LehalleNeuman2019,NeumanVoss2022,BankSonerVoss2017}, general propagators through infinite-dimensional stochastic control \citep{AbiJaberNeuman2025} and operator-resolvent methods \citep{AbiJaberNeumanTuschmann2024}, and power-law resilience with Gaussian signals on a finite horizon \citep{FordeSanchezSmith2022} --- the present treatment addresses the three frictions jointly on the whole line in the stationary regime, in closed transfer-function form, and adds the rate-response analysis of \S6.

\citet{AbiJaberNeuman2025} also treat general propagators with stochastic signals, by a different route. That approach solves the finite-horizon problem by infinite-dimensional stochastic control: the optimal rate is characterized as the unique solution of a linear stochastic Volterra equation of the second kind, with existence and uniqueness established for any Volterra propagator and any adapted signal, and closed forms recovered through kernel-specific reductions. The present approach works in the complementary stationary regime, using spectral factorization of the friction symbol: the projected-inverse identity \eqref{eq:pi-intro} expresses the adapted inverse of $Q$ as $Q_+^{-1}P_+Q_-^{-1}$ directly, without resolving a fixed-point equation. The two representations are equivalent on $[0,T]$ once the Gohberg--Krein factors of \S7 are substituted: the Volterra-equation resolvent series is the Neumann expansion of the inverse Volterra operator, which the factorization bypasses. In exchange for the stationary restriction, the factorization yields the trading filter \eqref{eq:filter}, the rate response \eqref{eq:response}, and the Markowitz--fractional interpolation of \S4--5, none of which are explicit in the Volterra-equation approach.

% =====================================================
\section{The Adapted Solution for General Signals}

\subsection{Setup}

Fix a filtered probability space satisfying the usual conditions. The signal $\alpha$ is adapted, mean-zero, with forecast curve
\[
\bar\alpha(t,s) \;=\; \begin{cases}\alpha_s, & s\le t,\\ \E_t[\alpha_s], & s > t,\end{cases}
\]
decaying as $s\to\infty$ (forecasts vanish at long horizons; this is also the admissibility condition $x\in L^2$ forced by the risk term). Our Fourier convention is $\hat f(\omega)=\int e^{i\omega t}f(t)\,dt$; causal kernels have transforms analytic in the upper half-plane. The symbol $Q(\omega)$ of \eqref{eq:symbol} is positive and even; we assume the Szeg\H{o} condition $\int|\log Q(\omega)|/(1+\omega^2)\,d\omega<\infty$ \citep{Krein1962}, satisfied by all kernels considered. Admissible controls are adapted processes of finite friction energy $\E\langle u, Qu\rangle<\infty$; for $\eta=0$ this class admits rates with a white-noise component (\S6), whose position $x$ remains in $L^2$. Then $Q$ admits a Wiener--Hopf factorization
\begin{equation}\label{eq:wh}
Q(\omega) = Q_-(\omega)\,Q_+(\omega), \qquad Q_- = \overline{Q_+},
\end{equation}
with $Q_+$ analytic and zero-free in the upper half-plane: $Q_+$ is causal with causal inverse, and $Q_- = Q_+^*$ is its anticausal adjoint. In computations it is convenient to clear the $\omega^{-2}$ pole of the risk term by factoring the \emph{entire} symbol
\begin{equation}\label{eq:N}
N(\omega) \;:=\; \omega^2 Q(\omega) \;=\; \eta\,\omega^2 + \gamma\,\hat G(\omega)\,\omega^2 + \lambda \;=\; N_-N_+,
\qquad Q_\pm(\omega) = \frac{N_\pm(\omega)}{\mp i\omega},
\end{equation}
with $N_+$ analytic and zero-free in the upper half-plane and $N_- = \overline{N_+}$ (the outer factor of $N$): the combined friction acts as a single propagator on the \emph{position}. The outer factor is given explicitly by the Szeg\H{o} formula
\begin{equation}\label{eq:szego}
N_+(\omega) \;=\; \exp\!\Bigl(\frac{1}{2\pi i}\int_{-\infty}^{\infty}\Bigl(\frac{1}{t-\omega}-\frac{t}{1+t^2}\Bigr)\log N(t)\,dt\Bigr), \qquad \operatorname{Im}\omega>0,
\end{equation}
unique up to a unimodular constant \citep{Krein1962}; since $N$ is even, the odd part of the integrand integrates to zero on the imaginary axis, and $N_+(i\theta)$ reduces to the real Poisson integral $\Phi(\theta)$ of \eqref{eq:phi}.

\subsection{The projected inverse and the general policy}

\begin{lemma}[Adapted projected inverse]\label{lem:pi}
For any factorization $q = Q_-Q_+$ as in \eqref{eq:wh}, on adapted processes
\begin{equation}\label{eq:pi}
(P_+\,Q\,P_+)^{-1} \;=\; Q_+^{-1}\,P_+\,Q_-^{-1}.
\end{equation}
\end{lemma}

The proof (Appendix A) is three lines of triangular algebra: causality makes $Q_+$ and its inverse preserve the adapted subspace, adjointness makes $Q_-$ and its inverse preserve the orthogonal complement, and the projection threads between them. The identity selects, among all square roots of $q$, the one whose whitening isometry $Q_-^{-1}$ leaves invariant the space of processes invisible to adapted strategies --- the sense in which the factorization is Cholesky with respect to the filtration.

\begin{theorem}[Optimal policy, general adapted signal]\label{thm:general}
The unique adapted optimum of \eqref{eq:objective} is
\begin{equation}\label{eq:policy}
u^\star \;=\; Q_+^{-1}\,P_+\,Q_-^{-1}\,\alpha,
\qquad\text{explicitly}\qquad
u^\star_t = \bigl(Q_+^{-1}\zeta\bigr)(t), \quad
\zeta_s = \bigl(Q_-^{-1}\,\bar\alpha(s,\cdot)\bigr)(s),
\end{equation}
where the second form uses the commutation of the deterministic operator $Q_-^{-1}$ with conditional expectation on the argument variable.
\end{theorem}

The policy has three steps. The anticausal operator $Q_-^{-1}$ applied to the signal requires its future path; the adapted projection replaces that future by the forecast curve, producing the $\F_s$-measurable \emph{whitened signal} $\zeta_s$; the causal operator $Q_+^{-1}$ then integrates $\zeta$ over the past to produce the rate. The architecture is Wiener--Kolmogorov prediction with the roles reversed: prediction whitens by the signal's own spectral factor to forecast it, trading whitens by the \emph{friction's} factor to monetize it.

\paragraph{Multiplier form and duality.} Dropping the adaptedness constraint, the maximization of \eqref{eq:objective} over $u$ is the Legendre transform of the friction energy $u\mapsto\tfrac12\E\langle u,Qu\rangle$ evaluated at the signal, with value $\tfrac12\E\langle\alpha,Q^{-1}\alpha\rangle$ attained at the anticipative optimum $u=Q^{-1}\alpha$ of \S1.1 (the value a trader would realize with the whole signal path in hand). Adaptedness restricts $u$ to a subspace, and the constrained transform carries a Lagrange multiplier. Substituting \eqref{eq:policy} into $Qu^\star$ writes the first-order condition \eqref{eq:foc} without the projection,
\begin{equation}\label{eq:multiplier}
Q\,u^\star \;=\; \alpha - \nu^\star, \qquad \nu^\star \;=\; Q_-\,(I-P_+)\,Q_-^{-1}\alpha, \qquad P_+\nu^\star = 0,
\end{equation}
so the adapted optimum solves the unconstrained first-order condition for the modified signal $\alpha-\nu^\star$, with $\nu^\star$ the multiplier of the adaptedness constraint. A multiplier with vanishing optional projection, $P_+\nu^\star=0$, is the defining property of the nonanticipativity multipliers of \citet{RockafellarWets1976}, established there for general multistage convex problems; \eqref{eq:multiplier} is their closed form for this problem. Every $\nu$ with $P_+\nu=0$ is orthogonal to the adapted controls, $\E\langle u,\nu\rangle=0$, so replacing $\alpha$ by $\alpha-\nu$ leaves the objective unchanged on adapted $u$ while relaxing it to the anticipative problem; the anticipative value $\tfrac12\E\langle\alpha-\nu,Q^{-1}(\alpha-\nu)\rangle$ therefore bounds the adapted value, and \eqref{eq:multiplier} attains the bound at $\nu^\star$. The constrained transform is thus a minimization over multipliers,
\begin{equation}\label{eq:dual}
\sup_{u\ \text{adapted}} \Bigl(\E\langle u,\alpha\rangle - \tfrac12\,\E\langle u,Qu\rangle\Bigr)
\;=\; \min_{P_+\nu = 0}\ \tfrac12\,\E\bigl\langle \alpha-\nu,\ Q^{-1}(\alpha-\nu)\bigr\rangle,
\end{equation}
the anticipative value of the signal with its non-adapted part removed, and the value forgone to adaptedness is the friction energy $\tfrac12\,\E\langle\nu^\star,Q^{-1}\nu^\star\rangle$ of the multiplier.

\paragraph{Mean-reverting signals.} A signal is \emph{mean-reverting} if its forecast curve factorizes into the current value and a deterministic lag profile,
\begin{equation}\label{eq:collapse}
\bar\alpha(t,s) \;=\; \alpha_t\, f(s-t) \qquad (s\ge t), \qquad f(0)=1,
\end{equation}
equivalently $\E_t[\alpha_{t+\tau}] = f(\tau)\,\alpha_t$. Consistency of conditional expectations forces the profile to be exponential: for $t_1<t_2<s$ the tower property gives
\[
\alpha_{t_1}f(s-t_1) = \E_{t_1}\!\bigl[\alpha_{t_2}f(s-t_2)\bigr] = \alpha_{t_1}\,f(t_2-t_1)\,f(s-t_2),
\]
so $f(a+b)=f(a)f(b)$ and $f(\tau)=e^{-\theta\tau}$, with $\theta>0$ (the mean-reversion rate) fixed by the forecast decay of \S2.1. Condition \eqref{eq:collapse} restricts conditional means only and is weaker than the Markov property: among stationary Gaussian signals it is the Ornstein--Uhlenbeck case (Doob's theorem), and it also holds for the non-Markov, non-Gaussian moving averages $\alpha_t = \int_{-\infty}^t e^{-\theta(t-r)}\sigma_r\,dW_r$ with adapted stochastic volatility $\sigma$.

Separability \eqref{eq:collapse} collapses the whitened signal to a multiple of the current value. In $\zeta_s = \bigl(Q_-^{-1}\bar\alpha(s,\cdot)\bigr)(s)$ the anticausal deterministic operator $Q_-^{-1}$ acts on the future profile $f(\cdot-s)$ while the $\F_s$-measurable factor $\alpha_s$ passes through, and stationarity of $Q_-^{-1}$ evaluates the shifted profile at lag zero,
\begin{equation}\label{eq:zeta-collapse}
\zeta_s \;=\; \rho\,\alpha_s, \qquad \rho \;=\; \bigl(Q_-^{-1}f\bigr)(0).
\end{equation}
The optimal policy is then a filter of the signal path, $u^\star = \rho\,Q_+^{-1}\alpha$: the whole forecast curve enters through the single number $\rho$. Vector mean reversion covers signals built from several timescales, where the profile is a matrix exponential: a read-out $\alpha_t = c^\top Y_t$ of a state with $\E_t[Y_{t+\tau}] = e^{-\Theta\tau}Y_t$ has $\bar\alpha(t,s) = c^\top e^{-\Theta(s-t)}Y_t$, whitened signal $\zeta_s = \rho^\top Y_s$ with $\rho^\top = \bigl(Q_-^{-1}[\,c^\top e^{-\Theta\cdot}\,]\bigr)(0)$, and a policy that filters the state path. For the OU signal $\rho = \theta/\Phi(\theta)$ with $\Phi(\theta)=N_+(i\theta)$ (\S3, Appendix~B); pure risk ($\Phi=\sqrt\lambda$, $Q_+^{-1}=\partial_t/\sqrt\lambda$) returns the Markowitz position $x^\star=\theta\alpha_t/\lambda$.

% =====================================================
\section{Stationary Signals: the Optimal Trading Filter}

Assume now that $\alpha$ is stationary, Gaussian, and purely non-deterministic (no perfectly predictable component; equivalently $\log S_\alpha$ is integrable), generating its own filtration, with spectral density $S_\alpha = |\varphi_+|^2$ and outer spectral factor $\varphi_+$ --- the causal, zero-free factor with $|\varphi_+|^2 = S_\alpha$. Then $\alpha = \varphi_+(D)\dot W$ is a causal moving average of its \emph{innovations} $\dot W$, the unit white noise generating the filtration \citep{Wiener1949}. Conditional expectation of a stationary linear functional of $\dot W$ truncates its moving-average kernel at lag zero, so $P_+$ acts on symbols as the Riesz projection $\Pi_+$ (truncation to strictly positive lags, with any lag-zero atom annihilated: the conditional expectation of a contemporaneous-innovation loading vanishes in the forward limit), and Theorem~\ref{thm:general} becomes a transfer function.

\begin{theorem}[Optimal trading filter]\label{thm:filter}
Under the standing assumptions,
\begin{equation}\label{eq:filter}
u^\star = \hat g(D)\,\dot W, \qquad
\hat g \;=\; Q_+^{-1}\,\Pi_+\!\bigl[Q_-^{-1}\varphi_+\bigr],
\end{equation}
and the attained value rate (expected P\&L per unit time) is
\begin{equation}\label{eq:value}
v \;=\; \frac{1}{4\pi}\,\bigl\|\Pi_+\bigl(Q_-^{-1}\varphi_+\bigr)\bigr\|_{L^2}^2 .
\end{equation}
\end{theorem}

The perfectly-anticipating benchmark replaces $\Pi_+$ by the identity, $v_{\rm ant} = \frac{1}{4\pi}\int S_\alpha/Q\,d\omega$, the value attainable if the whole signal path were known in advance. Its excess over $v$, the \emph{causality gap} $v_{\rm ant}-v$, is the value forgone by trading on current information alone; by \eqref{eq:value} it equals the $L^2$ energy of the anticausal part of the half-whitened signal $Q_-^{-1}\varphi_+$ (the spectral factor $\varphi_+$ with only the anticausal cost factor $Q_-^{-1}$ applied) --- per unit time, the friction energy of the nonanticipativity multiplier \eqref{eq:multiplier}. The position filter $\hat g_x = \hat g/(-i\omega)$ lies in $L^2$ for every kernel; the rate filter $\hat g$ lies in $L^2$ under the spectral-decay hypothesis of Appendix~B, and otherwise --- notably for $\eta=0$ kink kernels --- carries a lag-zero white-noise atom, the singular rate component analyzed in \S6.

For an Ornstein--Uhlenbeck signal --- mean reversion $\theta$, innovation variance $\sigma^2$, $\varphi_+ = \sigma/(\theta-i\omega)$ --- a pole cancellation collapses the projection (Appendix B), realizing \eqref{eq:zeta-collapse} in the frequency domain with $\rho=\theta/\Phi(\theta)$: $\Pi_+[Q_-^{-1}\varphi_+] = (\theta/\Phi(\theta))\,\varphi_+$ with
\begin{equation}\label{eq:phi}
\Phi(\theta) \;:=\; N_+(i\theta) \;=\; \exp\Bigl[\frac{\theta}{2\pi}\int\frac{\log N(t)}{\theta^2+t^2}\,dt\Bigr] \;>\;0,
\end{equation}
the causal factor of $N$ continued to the imaginary axis at $i\theta$, the point set by the signal's mean-reversion rate. The position transfer function with respect to the signal is then a single formula covering all three frictions:
\begin{equation}\label{eq:ou-filter}
\hat x^\star(\omega) \;=\; H(\omega)\,\hat\alpha(\omega), \qquad H(\omega) = \frac{\theta}{\Phi(\theta)\,N_+(\omega)},
\end{equation}
with rate filter $(-i\omega)H$, position response $X(\theta) := \E[x^\star_t\mid\alpha_t]/\alpha_t = \theta/\Phi(\theta)^2 > 0$, and value $v = \sigma^2\theta/4\Phi(\theta)^2 = \tfrac{\sigma^2}{4}X(\theta)$. Figure~\ref{fig:value} shows how this value depends on signal speed, and the cost of causality $v_{\rm ad}/v_{\rm ant}=\sin(\pi\beta/2)$.

\begin{figure}[t]
\centering
\includegraphics[width=\linewidth]{fig_value}
\caption{(a)~Value rate $v(\theta)$ against signal speed at fixed signal variance $\mathrm{Var}(\alpha)=1$ (normalized to $\theta=1$): faster signals are worth more under every friction, but transient impact makes the value grow only sub-linearly, $v\propto\theta^{1-\beta}$, against $\theta^2$ for pure risk. (Holding the innovation variance $\sigma^2$ fixed instead confounds speed with amplitude, since $\mathrm{Var}(\alpha)=\sigma^2/2\theta$.) (b)~The causality gap for an OU signal and power-law kernel, $v_{\rm ad}/v_{\rm ant}=\sin(\pi\beta/2)$: adaptedness is nearly free for short-memory impact ($\beta\to1$) and destroys the value for near-permanent impact ($\beta\to0$).}
\label{fig:value}
\end{figure}

% =====================================================
\section{Two Limits}

\paragraph{Markowitz: risk alone.} With $\eta=\gamma=0$, $N \equiv \lambda$ and the factorization is trivial: the filter is memoryless,
\[
x^\star_t \;=\; \frac{\E_t[\mu_t]}{\lambda}, \qquad \mu = -\dot\alpha \ \text{(the drift)},
\]
position equals conditionally expected return over risk \citep{Markowitz1952,Merton1971}. Memory is absent because the friction is local: the penalty prices the current position only, and each instant is a fresh mean--variance problem. The rate contains white noise --- positions jump to their targets --- which is the singularity any impact term regularizes. OU value rate: $v = \theta\sigma^2/4\lambda$.

\paragraph{Fractional derivative trading: power-law impact alone.} With $\eta=\lambda=0$ and $G = |t|^{-\beta}$, $\hat G = c_\beta|\omega|^{\beta-1}$, $c_\beta = 2\Gamma(1-\beta)\sin(\pi\beta/2)$, the factors are fractional integrals: $Q_\pm = (\gamma c_\beta)^{1/2} I_\pm^\nu$ with $\nu = (1-\beta)/2$, and inverses are Marchaud fractional derivatives $D_\pm^\nu$ \citep{SamkoKilbasMarichev1993}. Here $(I_+^\nu f)(t) = \Gamma(\nu)^{-1}\int_{-\infty}^t (t-s)^{\nu-1}f(s)\,ds$ is the Liouville fractional integral, $I_-^\nu$ its time reflection, and $D_\pm^\nu = (I_\pm^\nu)^{-1}$. Specializing the general policy \eqref{eq:policy},
\begin{equation}\label{eq:fractional}
u^\star_t \;=\; \frac{1}{\gamma c_\beta}\,(D_+^\nu\zeta)(t), \qquad
\zeta_s = \bigl(D_-^\nu\,\bar\alpha(s,\cdot)\bigr)(s):
\end{equation}
an anticausal fractional derivative of the forecast curve followed by a causal fractional derivative of the whitened signal, both of order $\nu\in(0.2,0.4)$ for empirical $\beta$; the constants $(\gamma c_\beta)^{-1/2}$ of the two factors are collected into the prefactor, so $\zeta$ here is $(\gamma c_\beta)^{1/2}$ times the whitened signal of Theorem~\ref{thm:general}. For an OU signal $\zeta_s = \theta^\nu\alpha_s$ and $u^\star = (\theta^\nu/\gamma c_\beta)D_+^\nu\alpha$, with value $v = \sigma^2\theta^{-\beta}/4\gamma c_\beta$. Memory is maximal: the Marchaud representation weights past signal increments by the slowly decaying kernel $(t-s)^{-1-\nu}$. On a grid $D_+^\nu$ is a Toeplitz operator whose Gr\"unwald--Letnikov weights inherit this tail (Figure~\ref{fig:gl}).

\begin{figure}[t]
\centering
\includegraphics[width=0.62\linewidth]{fig_gl_weights}
\caption{Gr\"unwald--Letnikov weights of the discrete fractional derivative $D_+^\nu$ (generalized binomial coefficients of $(1-z)^\nu$): the long-memory tail $|c_k|\sim k^{-(1+\nu)}$ is the discrete signature of the $(t-s)^{-1-\nu}$ Marchaud kernel.}
\label{fig:gl}
\end{figure}

The general filter interpolates between the memoryless and fractional poles across the crossover frequencies of $Q$: risk-dominated below $\omega_c = (\lambda/\gamma c_\beta)^{1/(1+\beta)}$, transient-impact-dominated in the middle, temporary-cost-dominated (myopic) above $\omega_* = (\gamma c_\beta/\eta)^{1/(1-\beta)}$. Figure~\ref{fig:filter} shows the resulting filter magnitude $|H(\omega)|\propto N(\omega)^{-1/2}$ across kernels.

\begin{figure}[t]
\centering
\includegraphics[width=\linewidth]{fig_trading_filter}
\caption{The trading filter in the frequency domain, $|H(\omega)|\propto N(\omega)^{-1/2}$. (a)~Power-law impact ($\eta=\lambda=0$) gives a scale-free fractional filter of slope $-(1+\beta)/2$. (b)~The friction family ($\gamma=1$, $\kappa=2$, $\lambda=0.5$, and $\eta=0.3$ where present): inventory risk sets a low-frequency plateau, exponential transient impact a resilience shelf, temporary cost a high-frequency rolloff, and power-law impact a fractional tail, with crossover frequencies $\omega_c,\omega_*$.}
\label{fig:filter}
\end{figure}

% =====================================================
\section{Explicit Filters and Earlier Solutions}

Rational symbols factor in closed form and produce filters with finitely many exponential moving averages; this section computes them and locates the known solutions of the literature inside the family. The smoothing rates are the poles of $N_+$, and inventory risk sets their scale (Figure~\ref{fig:structure}).

\begin{figure}[t]
\centering
\includegraphics[width=\linewidth]{fig_filter_structure}
\caption{The optimal policy across kernels and parameters. (a)~For the temporary-plus-resilient kernel ($\eta=0.5$, $\gamma=1$, $\kappa=2$), the two moving-average rates $b_1,b_2$ against risk aversion: the slow rate $b_1\to0$ as $\lambda\to0$ (the position becomes non-stationary without risk) and both grow with $\lambda$. (b)~The position impulse response $g_x(\tau)$ by kernel ($\theta=1$, normalized at $\tau=0.3$): the exponential and two-EMA kernels have finite memory (exponential cutoff), the power-law kernel a heavy tail, and the Markowitz limit is an instantaneous $\delta$. (c)~The power-law policy across $\beta$: a scale-free kernel of slope $(\beta-1)/2$.}
\label{fig:structure}
\end{figure}

\subsection{Exponential propagator with risk: signal plus one moving average}

For $G = e^{-\kappa|t|}$, $\eta=0$:
\begin{equation}\label{eq:exp-factor}
N(\omega) = \frac{A\omega^2 + \lambda\kappa^2}{\kappa^2+\omega^2}, \qquad
N_+ = \sqrt A\,\frac{m-i\omega}{\kappa-i\omega}, \qquad
A = 2\kappa\gamma+\lambda,\quad m = \kappa\sqrt{\lambda/A},
\end{equation}
and \eqref{eq:ou-filter} splits into
\begin{equation}\label{eq:ema}
x^\star_t \;=\; \frac{\theta(\kappa+\theta)}{A(m+\theta)}\Bigl[\alpha_t + (\kappa-m)\!\int_{-\infty}^t e^{-m(t-s)}\alpha_s\,ds\Bigr]:
\end{equation}
the current signal plus one EMA at rate $m$, which travels from $0$ to $\kappa$ as risk aversion grows. Differentiating at $\lambda\to0$ gives the pure-impact rate $u^\star = \frac{\kappa+\theta}{2\kappa\gamma}(\dot\alpha+\kappa\alpha)$; at $\lambda\to\infty$ the EMA weight vanishes and the Markowitz filter returns.

\subsection{Temporary cost with risk: the aim portfolio recovered}

With $\gamma=0$, $N = \eta\omega^2+\lambda$, $N_+ = \sqrt\eta\,(a-i\omega)$, $a = \sqrt{\lambda/\eta}$: one EMA, and differentiation puts the filter in partial-adjustment form
\begin{equation}\label{eq:gp}
u^\star_t = a\,\bigl(\mathrm{aim}_t - x^\star_t\bigr), \qquad
\mathrm{aim}_t = \frac{a}{a+\theta}\cdot\frac{\theta\alpha_t}{\lambda}:
\end{equation}
trade at speed $\sqrt{\lambda/\eta}$ toward an aim that discounts the Markowitz position by the persistence factor $a/(a+\theta)$. This is the continuous-time solution of \citet{GarleanuPedersen2016} (discrete time: \citealp{GarleanuPedersen2013}): a diagonal friction produces exactly one pole. Figure~\ref{fig:aim} illustrates the partial adjustment.

\begin{figure}[t]
\centering
\includegraphics[width=\linewidth]{fig_aim_portfolio}
\caption{%
(a)~Optimal positions on the same OU signal path ($\theta=\lambda=1$, $\mathrm{Var}(\alpha)=1$, $\eta=0.5$ throughout), with the Markowitz target $\theta\alpha/\lambda$ (grey dashed) as reference. All three regimes damp positions relative to Markowitz; more persistent impact means a smaller optimal position (standard deviations: $0.375$, $0.328$, $0.269$ vs $0.87$ for Markowitz). The aim portfolio holds the largest of the three positions because temporary cost alone imposes no penalty on position size; adding exponential or power-law transient impact discourages holding large inventories. (b)~Rate step response: the analytical formula for the optimal trading rate per unit of a permanent step in $\alpha$, from the transfer functions of \S\S4--5, shown on a log lag axis. The aim portfolio (one EMA, decay rate $a=\sqrt{\lambda/\eta}$) and the NV stationary solution (two EMAs with rates $b_1{<}a{<}b_2$) both decay exponentially; the slow NV component $b_1{<}a$ makes the NV response more persistent than the aim portfolio. The power-law regime decays algebraically, $\propto \tau^{-(1-\beta)/2}$, the time-domain signature of fractional memory.}
\label{fig:aim}
\end{figure}

\begin{figure}[t]
\centering
\includegraphics[width=\linewidth]{fig_interpolation}
\caption{%
Interpolation between the fractional-derivative and aim-portfolio limits,
fixed power-law kernel ($\beta=0.5$, $\gamma=\lambda=\theta=1$), varying temporary-cost strength $\eta$.
(a)~Adapted positions (stationary Szeg\H{o} filter applied to the same OU signal path as Figure~\ref{fig:aim}), with the Markowitz target $\theta\alpha/\lambda$ (grey dashed) for reference.
All positions are damped relative to Markowitz and to each other: adding temporary cost ($\eta$) shrinks the optimal position because each unit of trading rate is penalized directly, forcing more conservative inventory targets.
The two small-$\eta$ curves (green) are nearly indistinguishable because the power-law transient cost dominates both at the signal bandwidth; the divergence becomes pronounced only for $\eta\gtrsim1$ where the temporary-cost crossover frequency $\xi_*=(\gamma c_\beta/\eta)^{2}$ falls below $\theta$.
(b)~Rate step response $g_x(\tau)$ --- the optimal trading rate per unit of a permanent step in $\alpha$ --- computed via the Szeg\H{o} outer factor of the combined friction symbol $N(\omega)=\eta\omega^2+\gamma c_\beta|\omega|^{1+\beta}+\lambda$, normalised at $\tau=0.3$, on a log lag axis. The dotted curve is the pure fractional limit ($\eta\to0$, $\propto\tau^{-(1-\beta)/2}$) and the dashed curve is the GP exponential limit ($\gamma\to0$, $\propto e^{-at}$, $a=\sqrt{\lambda/\eta}$). The five coloured curves sit between the two analytical limits and converge monotonically from the power-law tail (small $\eta$) to the exponential decay (large $\eta$) as temporary cost grows.
}
\label{fig:interp}
\end{figure}

\subsection{Temporary plus resilient impact: Neuman--Vo\ss{} in filter form}

\citet{NeumanVoss2022} treat temporary cost with exponential-resilience transient impact on a finite horizon, characterizing the optimum through four coupled linear FBSDEs solved by matrix exponentials, in feedback form on the inventory and the impact state. The stationary counterpart (infinite horizon, terminal penalty dropped) is explicit; for the exponential kernel the three-friction problem is also solved in feedback form by \citet[Prop.~2]{GarleanuPedersen2016}. With all three frictions the symbol is rational with biquadratic numerator,
\begin{equation}\label{eq:nv-factor}
N_+ = \sqrt\eta\,\frac{(b_1-i\omega)(b_2-i\omega)}{\kappa-i\omega}, \qquad
b_1^2b_2^2 = \frac{\lambda\kappa^2}{\eta}, \qquad
b_1^2+b_2^2 = \kappa^2 + \frac{2\kappa\gamma+\lambda}{\eta},
\end{equation}
and partial fractions turn \eqref{eq:ou-filter} into a two-EMA filter,
\begin{equation}\label{eq:nv-filter}
x^\star_t \;=\; \frac{\theta}{\Phi(\theta)\sqrt\eta}\Bigl[w_1\!\int_{-\infty}^t\! e^{-b_1(t-s)}\alpha_s ds + w_2\!\int_{-\infty}^t\! e^{-b_2(t-s)}\alpha_s ds\Bigr],
\qquad w_i = \frac{\kappa-b_i}{b_j-b_i}.
\end{equation}
The FBSDE feedback data reduce, in the stationary regime, to two smoothing rates and two weights --- roots and residues of the symbol; the finite-horizon feedback gains converge to these stationary values in the interior of the horizon (Figure~\ref{fig:nv}), with the Riccati closed-loop poles equal to the moving-average rates $b_1,b_2$. Both degenerations are exact: at $\gamma=0$ the roots are $\{a,\kappa\}$ and the $\kappa$-pole cancels the kernel zero, collapsing to \eqref{eq:gp}; as $\eta\to0$, $b_1\to m$ and $\sqrt\eta(b_2-i\omega)\to\sqrt A$, recovering \eqref{eq:ema}. The pattern generalizes: a rational friction with $n$ poles yields a filter with at most $n$ EMAs; the power-law kernel is a scale-free superposition of exponentials, $|t|^{-\beta} = \Gamma(\beta)^{-1}\int_0^\infty e^{-r|t|}r^{\beta-1}\,dr$, and yields the fractional filter with a continuum of EMA rates.

\begin{figure}[t]
\centering
\includegraphics[width=\linewidth]{fig_nv_vs_stationary}
\caption{Finite-horizon (Neuman--Vo\ss) solution against the stationary filter; $\eta=0.5$, $\gamma=1$, $\kappa=2$, $\lambda=1$, $T=20$. (a)~The finite-horizon feedback gains $K(t)$ on inventory and impact state sit on the stationary constants $K_\infty$ (dashed) through the interior, departing only in the terminal boundary layer (backward Riccati). (b)~Optimal position for $\alpha_t=\sin\omega_0 t$: the finite-horizon rule (red) tracks the whole-line stationary filter (black) in the interior and deviates in the shaded boundary layers --- start-up from zero inventory and terminal unwind. The Riccati closed-loop poles equal the filter's moving-average rates $b_1,b_2$.}
\label{fig:nv}
\end{figure}

% =====================================================
\section{Impact Surfing and the Rate Response}

\subsection{The response formula}

For the OU signal, define the flow response $R(\theta) = \lim_{\epsilon\downarrow0}\E[u^\star_{t+\epsilon}\mid\alpha_t]/\alpha_t$: the average direction of the trade executed just after $t$ by a trader holding the current signal value. The forward limit handles the distribution-valued rate. For $\eta=0$ kink kernels $u^\star_t$ carries a contemporaneous white-noise atom $a_0\dot W_t$ (Appendix~C) correlated with $\alpha_t$, so conditioning at the same instant would fold it in; evaluating an instant later, where the innovations in $(t,t+\epsilon]$ are independent of $\alpha_t$, drops the atom and leaves the execution-relevant response. The position response $X(\theta)$ below needs no limit, the position filter $\hat g_x$ being square-integrable.

\begin{proposition}\label{prop:response}
For rational or power-law symbols,
\begin{equation}\label{eq:response}
R(\theta) \;=\; \frac{\theta^2}{\Phi(\theta)}\Bigl[\frac{1}{\Phi(\theta)} - 2c_1\Bigr],
\qquad c_1 = \lim_{|\omega|\to\infty}\frac{1}{N_+(\omega)},
\end{equation}
while the position response $X(\theta) = \theta/\Phi(\theta)^2$ is positive for every kernel: positions never flip; flows flip exactly when $2c_1\Phi(\theta) > 1$.
\end{proposition}

The proof is in Appendix~C. The constant $c_1$ is the innovation atom of the rate filter: the weight the optimal rate places on the contemporaneous innovation, present exactly when singular rates carry finite cost. With $\eta = 0$, $c_1 = [\lambda - 2\gamma G'(0^+)]^{-1/2}$ for a \emph{kink kernel} --- finite slope $G'(0^+)$ at the origin, as for $e^{-\kappa|t|}$ --- and $c_1 = 0$ for the power-law cusp ($G'(0^+)=-\infty$); any temporary cost $\eta>0$ also sets $c_1 = 0$.

\subsection{The impact-surfing phase diagram}

\emph{Exponential propagator, no temporary cost.} $\Phi = \sqrt A(m+\theta)/(\kappa+\theta)$ gives a single flip at
\begin{equation}\label{eq:threshold}
\theta^* \;=\; \kappa - 2m \;=\; \kappa\Bigl[1 - 2\sqrt{\lambda/(2\kappa\gamma+\lambda)}\Bigr].
\end{equation}
At $\lambda = 0$ the threshold is $\theta^* = \kappa$: when the signal mean-reverts faster than the impact decays, a signal-following trade would leave an impact residual that outlives the signal, and the optimal flow instead trades against the signal to ride that residual --- \emph{impact surfing}. Inventory risk lowers the threshold, and for $\lambda \ge 2\kappa\gamma/3$ the optimal flow surfs at every signal speed: the position \eqref{eq:ema} stays aligned with the signal, but its mean reversion dominates the conditional flow. The pure-risk limit $R = -\theta^2/\lambda$, $X = \theta/\lambda$ is the extreme case: a Markowitz position is always conditionally exiting. Figure~\ref{fig:surf} shows the response and the resulting phase diagram.

\begin{figure}[t]
\centering
\includegraphics[width=\linewidth]{fig_impact_surfing}
\caption{Impact surfing. (a)~The flow response $R(\theta)$ for the exponential kernel ($\gamma=1$, $\kappa=2$) changes sign at $\theta^*=\kappa-2m$ as inventory risk $\lambda$ grows; the power-law kernel never flips. (b)~Phase diagram in $(\theta,\lambda)$: the optimum follows the signal ($R>0$) below $\theta^*(\lambda)$ and surfs its own transient impact ($R<0$) above, surfing at every signal speed once $\lambda\ge 2\kappa\gamma/3$.}
\label{fig:surf}
\end{figure}

\emph{Power-law propagator.} $c_1 = 0$ at every $\lambda$: flow and position are aligned at all signal speeds and risk aversions. Impact surfing is confined to kink kernels; neither inventory risk nor forecast speed induces it under a power-law cusp.

\emph{Any temporary cost.} $c_1 = 0$: the flow partial-adjusts toward an aligned aim and $R = \theta^2/\Phi^2 > 0$. Impact surfing is carried entirely by the singular (block-trade) component of the optimal policy, which temporary cost removes.

\emph{Impact surfing and price manipulation.} Impact surfing profits from the trader's own transient-impact wake rather than from the signal, so it borders on the strategies that impact models are designed to exclude. The positive-definiteness $\hat G\ge0$ of \eqref{eq:symbol} rules out static round-trip arbitrage and, for the exponential and power-law kernels used here, transaction-triggered price manipulation \citep{Gatheral2010,AlfonsiSchiedSlynko2012}: surfing is a cost-optimal response, not an arbitrage. It is nonetheless the regime in which the optimum most exploits self-impact, and the frictions omitted here --- bid--ask spread, position and short-sale limits, order discreteness --- would temper it, so the surfing prediction should be read within the model class rather than as a standalone recommendation.

\subsection{Numerical verification}

Discretized adapted optima reproduce the closed forms. On a uniform grid of spacing $dt$, the cost form is the matrix $\eta I + \gamma\,dt\,G + \lambda\,dt^2 L^\top\!L$, where $G$ is the sampled kernel and $L$ the lower-triangular cumulative-sum matrix mapping rates to positions; the adapted optimum is computed by reverse-order Cholesky (causal factor on the right) and the forward responses by lag-one regression on the signal (grids: $dt=0.01$ for rows~1 and~4, $dt=0.04$ for rows~2--3, $dt=0.02$ for row~5):

\begin{center}
\begin{tabular}{lccccc}
\toprule
Kernel & $\eta$ & $\lambda$ & $\theta$ & $R$ (formula / discrete) & $X$ (formula / discrete)\\
\midrule
$e^{-2|t|}$, $\gamma{=}1$ & 0 & $0.5$ & $1.5$ & $-0.311$ / $-0.309$ & $0.870$ / $0.860$\\
$e^{-2|t|}$, $\gamma{=}1$ & 0 & $4$ & $0.5$ & $-0.028$ / $-0.028$ & $0.107$ / $0.104$\\
none, $\gamma{=}0$ & 0 & $1$ & $0.7$ & $-0.490$ / $-0.477$ & $0.700$ / $0.671$\\
$|t|^{-1/2}$, $\gamma{=}1$ & 0 & $1$ & $2$ & $+0.364$ / $+0.325$ & $0.182$ / $0.190$\\
$e^{-2|t|}$, $\gamma{=}1$ & $0.5$ & $1$ & $1$ & $+0.264$ / $+0.241$ & $0.264$ / $0.271$\\
\bottomrule
\end{tabular}
\end{center}

Signs match everywhere, including the always-surfing regime ($\lambda = 4 > 2\kappa\gamma/3$) and the sign restoration by temporary cost (last row). The exponential-family entries converge to the formulas as $dt\to0$ (first order; row~1 reaches $0.5\%$ in $R$ at $dt=0.01$), and the residual power-law discrepancy is comparable to the quadrature bias of the singular kernel measured at the analytically known $\lambda=0$ point on the same grid.

% =====================================================
\section{Boundary Effects and the Gohberg--Krein Factorization}

On a finite horizon $[0,T]$ --- liquidation, daily sessions --- translation invariance is lost and Fourier factorization no longer applies. The Wiener--Hopf factors are replaced by the triangular factorization of positive operators relative to the nest of projections $\{P_{[0,t]}\}$ (increasing in $t$) \citep{GohbergKrein1970,Arveson1975}. For the operators arising here --- $\eta I$ plus a continuous-kernel integral operator, and the explicitly exhibited power-law factor below --- the cost operator restricted to $[0,T]$, written $G_T$, factors as $G_T = C_-C_+$ with $C_+$ a causal Volterra operator and $C_- = C_+^*$ \citep{GohbergKrein1970,Arveson1975}. Lemma~\ref{lem:pi} and Theorem~\ref{thm:general} apply with these factors:
\begin{equation}\label{eq:finiteT}
u^\star \;=\; C_+^{-1}\,P_+\,C_-^{-1}\,\alpha^{\rm eff},
\end{equation}
where $\alpha^{\rm eff} = \alpha + \sum_k \mu_k\psi_k$ adjoins a multiplier $\mu_k$ for each position constraint, determined by the constraint equations (a process-valued multiplier for the pathwise liquidation constraint $x_T = 0$); position constraints thus modify the signal in the same way the adaptedness constraint does in \eqref{eq:multiplier}, each multiplier confined to its constraint's annihilator. The identity \eqref{eq:pi} requires the causal factor \emph{on the right}, which pins the factor's anchoring. For the power-law kernel ($\eta=\lambda=0$) the explicit factor is the terminal-anchored Volterra operator
\begin{equation}\label{eq:gk-kernel}
(C_+f)(t) = (\gamma c_\beta)^{1/2}\!\int_0^t \Bigl(\frac{T-s}{T-t}\Bigr)^{\!\nu}\frac{(t-s)^{\nu-1}}{\Gamma(\nu)}\,f(s)\,ds,
\end{equation}
verified by direct kernel integration ($C_-C_+ = G_T$ exactly, including the constant). The time reflection of \eqref{eq:gk-kernel} is the left-anchored outer factor $G_T = \mathcal T\mathcal T^*$ used by \citet[pp.~590--591]{FordeSanchezSmith2022} to invert $G_T$ \emph{in full}, where factor order is immaterial; the adapted projection between the factors is what forces the terminal anchoring.

\emph{Recovering earlier solutions.} The first-order condition of the finite-horizon problem is the stochastic Fredholm equation that \citet{FordeSanchezSmith2022} solve coefficient-wise on the Wiener chaos for Gaussian Volterra signals (their Theorem 2.2); formula \eqref{eq:finiteT} is its solution operator, and expanding \eqref{eq:finiteT} on the chaos reproduces their Fredholm systems, with the generalized Abel inversion of \citet{ChakrabartiGeorge1994} implementing the full inverse $G_T^{-1}$ and weighted Riemann--Liouville operators implementing the individual factors. For a constant signal --- pure liquidation --- the projection is trivial and \eqref{eq:finiteT} reduces to $G_T^{-1}\mathbf 1$ up to the terminal-inventory multiplier, the U-shaped profile $u^\star(t) \propto [t(T-t)]^{(\beta-1)/2}$ of \citet[Ex.~2.30]{GatheralSchiedSlynko2012}.

\emph{Boundary layers.} In the interior the terminal-anchored weights $((T-s)/(T-t))^\nu$ tend to one and \eqref{eq:gk-kernel} tends to the whole-line fractional factor, so the finite-horizon policy converges to the stationary filter away from the boundaries. The rate of convergence depends on the kernel.

\begin{proposition}[Boundary-layer decay]\label{prop:boundary}
Let $u^{\star,T}$ be the finite-horizon optimum on $[0,T]$ and $u^\star$ the whole-line stationary optimum for the same bounded signal, and $d(t)=\min(t,T-t)$ the distance to the nearest endpoint. Away from the boundaries the correction decays with $d(t)$:
\begin{equation}\label{eq:bdry-decay}
\bigl|u^{\star,T}_t - u^\star_t\bigr| \;\le\;
\begin{cases}
C(\beta)\,\|\alpha\|_\infty\,d(t)^{-\nu}, & \text{power-law kernel},\ \nu=\tfrac{1-\beta}{2},\\[3pt]
C\,\|\alpha\|_\infty\,e^{-b_1 d(t)}, & \text{rational symbol, slowest causal-factor rate } b_1>0.
\end{cases}
\end{equation}
\end{proposition}

The finite-horizon policy \eqref{eq:finiteT} differs from the stationary one through the terminal-anchored weights and the truncation of the factor kernels to $[0,T]$. For the power-law factor \eqref{eq:gk-kernel} the truncation cuts the Marchaud tail at distance $d(t)$, contributing $d(t)^{-\nu}$; the weight deviation $((T-s)/(T-t))^\nu = 1+O(d(t)^{-1})$ is subdominant. For a rational symbol the causal factor $N_+^{-1}$ has kernel decaying as $e^{-b_1\tau}$ at the slowest zero $b_1$ of $N_+$, so truncation at distance $d(t)$ costs $e^{-b_1 d(t)}$; equivalently the finite-horizon feedback gains relax to their stationary values at rate $2b_1$ (Figure~\ref{fig:nv}). Boundary effects are thus confined to layers of width $\sim 1/b_1$ for rational kernels and set by the impact memory for power laws, large only for the small $\nu$ of empirical exponents --- a quantitative expression of how long-memory impact makes horizons matter.

% =====================================================
\section{Concluding Remarks}

The Wiener--Hopf method turns the adapted gain--risk--cost problem into algebra: factor the friction symbol into causal and anticausal parts, thread the optional projection between them, and read the optimal policy off the factors. The classical portfolio and execution rules are the method's degenerate cases --- Markowitz when the friction is local in the position, the aim portfolio when it is local in the rate, fractional-derivative trading when it is scale-free --- and the filter language makes the interpolation explicit, down to the number of moving averages (the poles of the symbol) and the impact-surfing threshold (the zero of the rate response).

The closed forms hold under several restrictions. Explicit filters require stationary Gaussian signals generating their own filtration, and the OU formulas use the mean-reversion condition \eqref{eq:collapse}; the general policy \eqref{eq:policy} holds for any adapted signal but leaves the forecast curve to the user. Costs are quadratic and the kernel symmetric: spread, tick effects, and nonlinear impact are outside the model. For the pure power-law kernel the factors are unbounded operators, and the formulas hold on a dense domain fixed by a spectral-decay hypothesis; any $\eta>0$ or $\lambda>0$ restores bounded inverses. On finite horizons the theory is complete only up to the boundary terms attached to constraints, for which \eqref{eq:finiteT} gives the operator form but not yet sharp constants.

% =====================================================
\appendix

\section{Proof of Lemma \ref{lem:pi} and Theorem \ref{thm:general}}
\emph{Lemma \ref{lem:pi}.} Causality of $Q_+$ gives $P_+^\perp Q_+P_+ = 0$, so $Q_+$ and its causal inverse preserve the adapted subspace; the adjoint identity $P_+Q_-P_+^\perp = 0$ shows $Q_-$ and $Q_-^{-1}$ preserve its orthogonal complement. For adapted $v$, split $Q_-v = P_+Q_-v + P_+^\perp Q_-v$; applying $Q_-^{-1}$, the image of the second term stays in the complement and is annihilated by $P_+$, so $P_+Q_-^{-1}P_+Q_-v = P_+v = v$. Hence for adapted $u$,
\[
Q_+^{-1}P_+Q_-^{-1}\,(P_+QP_+)\,u = Q_+^{-1}\bigl[P_+Q_-^{-1}P_+Q_-\bigr](Q_+u) = Q_+^{-1}Q_+u = u .
\]
The reverse identity $(P_+QP_+)(Q_+^{-1}P_+Q_-^{-1})\alpha = \alpha$ for adapted $\alpha$ follows the same way (using that $Q_-$ preserves the orthogonal complement), or from positivity of $P_+QP_+$ on the adapted subspace, for which a left inverse is two-sided. For the pure power-law kernel the factors are unbounded and the computation holds on the dense domain fixed by the spectral-decay hypothesis (Appendix~B); $\eta>0$ or $\lambda>0$ gives bounded inverses.

\emph{Theorem \ref{thm:general}.} Strict convexity ($q>0$) makes the projected first-order condition \eqref{eq:foc} necessary and sufficient with a unique solution, which Lemma~\ref{lem:pi} inverts: $u^\star = Q_+^{-1}P_+Q_-^{-1}\alpha$. The forecast-curve form follows from conditional Fubini: $Q_-^{-1}$ is a deterministic anticausal kernel, so $\E_s\bigl[(Q_-^{-1}\alpha)(s)\bigr]$ integrates that kernel against $\E_s[\alpha_r] = \bar\alpha(s,r)$, $r\ge s$.

\section{Proof of Theorem \ref{thm:filter} and the OU formulas}
The process $Q_-^{-1}\alpha$ is the stationary filter of $\dot W$ with symbol $h := Q_-^{-1}\varphi_+$; conditional expectation at $s$ annihilates innovations after $s$, so $P_+$ truncates the moving-average kernel at lag zero, acting as $\Pi_+$ on the symbol; atoms at lag zero load on the contemporaneous innovation and are likewise annihilated. Composing with $Q_+^{-1}$ gives \eqref{eq:filter}. Membership $\hat g\in L^2$: writing $h = N_-^{-1}(i\omega)\varphi_+$, the hypothesis $\int(1+\omega^2)\,S_\alpha(\omega)/Q(\omega)\,d\omega<\infty$ places $h$ in a Sobolev class on which the half-line indicator multiplies boundedly \citep[Ch.~1]{LionsMagenes1972}. For the value, the first-order condition and adaptedness give $v = \tfrac12\E[u^\star_t\alpha_t]$, and by the It\^o isometry and $\overline{\varphi_+} = Q_+\overline h$ on the real axis, $\E[u^\star\alpha] = \tfrac{1}{2\pi}\langle\Pi_+h,h\rangle = \tfrac{1}{2\pi}\|\Pi_+h\|^2$. For the OU signal, $h = \sigma i\omega/(N_-(\omega)(\theta-i\omega))$ and
\[
h - \frac{\theta}{\Phi(\theta)}\varphi_+ = -\sigma N_-^{-1} + \sigma\theta\Bigl[\frac{1}{N_-(\omega)} - \frac{1}{N_-(-i\theta)}\Bigr]\frac{1}{\theta-i\omega}
\]
is anticausal-plus-atom: the bracket's pole at $\omega=-i\theta$ cancels, $1/N_-$ is analytic in the lower half-plane, and $N_-(-i\theta) = N_+(i\theta) = \Phi(\theta)$ by conjugate symmetry (real and positive since $N$ is even and positive). Hence $\Pi_+h = (\theta/\Phi)\varphi_+$, giving \eqref{eq:ou-filter}, and $\|\varphi_+\|^2 = \pi\sigma^2/\theta$ gives $v = \sigma^2\theta/4\Phi^2$.

\section{Proof of Proposition \ref{prop:response}}
The position filter $\hat g_x = (\theta/\Phi)N_+^{-1}\varphi_+$ vanishes at high frequency for every kernel considered, so $X(\theta)$ is the plain Laplace evaluation $\theta/\Phi(\theta)^2 > 0$. The rate filter $\hat g = (-i\omega)(\theta/\Phi)N_+^{-1}\varphi_+$ tends to $a_0 = \theta\sigma c_1/\Phi$; splitting $g = a_0\delta + g_{\rm reg}$, the atom is excluded from the forward limit (the innovations in $(t,t+\epsilon]$ are independent of $\alpha_t$), and dominated convergence gives $R = (2\theta/\sigma^2)\,\sigma[\hat g(i\theta) - a_0] = \theta^2\Phi^{-1}[\Phi^{-1}-2c_1]$. With $\eta=0$ and a kink kernel, $\hat G(\omega)\omega^2\to-2G'(0^+)$ so $N(\infty) = \lambda - 2\gamma G'(0^+) = c_1^{-2}$; power-law cusps and any $\eta>0$ make $N_+$ grow, $c_1=0$. Consistency: $\lambda\to0$ exponential returns $(\kappa^2-\theta^2)/2\kappa\gamma$; $\gamma\to0$ returns $-\theta^2/\lambda$, matching $x = \theta\alpha/\lambda$ directly.

\paragraph{Acknowledgements.} The author thanks colleagues at CFM for discussions. Drafting, derivation checking, and numerical verification were assisted by Claude Opus via the Feynman research harness; all mathematical claims were verified by the author.

\bibliographystyle{plainnat}
\bibliography{optimal-trading-filters}

\end{document}
